\documentclass[10pt,a4paper]{article}

%double si on veut une version prof et une version élève. simple si on veut une seule version.
\newcommand{\buildmode}{simple}

\InputIfFileExists{version.tex}{}{}
% Version (définie par le script)
\providecommand{\version}{eleve}

\usepackage{feuilleinterro}

%%%à modifier à chaque fois

\renewcommand{\niveau}{2nde}
\renewcommand{\chapitre}{Équations, solution d'inéquations}
\renewcommand{\numchap}{0.3}
\renewcommand{\theme}{Nombres et calcul}

%%%titre "CORRECTION" au lieu de "INTERROGATION", sans les champs Nom/Prénom
\renewcommand{\interroversion}[1]{
    \noindent\textbf{\niveau}

    \vspace{-0.5cm}
    \begin{center}
        \textbf{\large CORRECTION -- \chapitre}
    \end{center}

    \vspace{0.1cm}
}

%%%axe gradué (identique au sujet) avec la solution placée en rouge
%%% #1 : abscisse (valeur numérique) du point à placer
%%% #2 : étiquette (code LaTeX) affichée au-dessus du point
\newcommand{\axegraduesolution}[2]{
\begin{center}
\begin{tikzpicture}[scale=1.5]
\draw[->] (-5.5,0)--(5.5,0);
\foreach \x in {-5,..., 4}{
    \draw (\x,0.1)--(\x,-0.1) node[below] {\(\x\)};
    \draw (\x+0.5,0.07)--(\x+0.5,-0.07) ;
    \foreach \minx in {0.25, 0.5, 0.75}{
        \draw (\x+\minx, 0.05)--(\x+\minx, -0.05);
    }
}
    \draw (5,0.1)--(5,-0.1) node[below] {\(5\)};
    \filldraw[red] (#1,0) circle (1.5pt) node[above]{\scriptsize \(#2\)};
\end{tikzpicture}
\end{center}
}

\begin{document}

% =========================================================
% PAGE 1 : VERSION A
% =========================================================

\interroversion{A}

\textbf{Exercice 1 -- Écriture décimale de fractions \hfill /2}

\begin{enumerate}
\item \(\dfrac{7}{4}=\textcolor{blue}{1{,}75}\)

\item \(-\dfrac{3}{5}=\textcolor{blue}{-0{,}6}\)
\end{enumerate}

\textbf{Exercice 2 -- Solution d'une équation \hfill /4}

Testons \(x=3\) dans \((E_1)\) :
\[
2\times3+1=6+1=7.
\]
On retrouve bien le membre de droite, donc \(\textcolor{blue}{3 \text{ est solution de } (E_1)}\).

Testons \(x=3\) dans \((E_2)\) :
\[
3\times3-2=9-2=7\neq10.
\]
Donc \(\textcolor{blue}{3 \text{ n'est pas solution de } (E_2)}\).

\textbf{Exercice 3 -- Résolution d'équation \hfill /4}

\[
\begin{aligned}
2x+1&=4\\
2x&=4-1\\
2x&=3\\
x&=\dfrac{3}{2}.
\end{aligned}
\]

L'ensemble des solutions est \(\textcolor{blue}{S=\left\{\dfrac{3}{2}\right\}}\).

\axegraduesolution{1.5}{x=\frac{3}{2}}

\textbf{Exercice 4 -- Intervalles (bonus) \hfill /1}

La borne \(-2\) est incluse (\(\leqslant\)) et la borne \(5\) est exclue (\(<\)), donc l'intervalle recherché est
\[
\textcolor{blue}{\left[-2\ ;\ 5\right[}.
\]

\newpage

% =========================================================
% PAGE 2 : VERSION B
% =========================================================

\interroversion{B}

\textbf{Exercice 1 -- Écriture décimale de fractions \hfill /2}

\begin{enumerate}
\item \(\dfrac{9}{2}=\textcolor{blue}{4{,}5}\)

\item \(-\dfrac{11}{4}=\textcolor{blue}{-2{,}75}\)
\end{enumerate}

\textbf{Exercice 2 -- Solution d'une équation \hfill /4}

Testons \(x=4\) dans \((E_1)\) :
\[
3\times4-5=12-5=7.
\]
On retrouve bien le membre de droite, donc \(\textcolor{blue}{4 \text{ est solution de } (E_1)}\).

Testons \(x=4\) dans \((E_2)\) :
\[
2\times4+3=8+3=11\neq13.
\]
Donc \(\textcolor{blue}{4 \text{ n'est pas solution de } (E_2)}\).

\textbf{Exercice 3 -- Résolution d'équation \hfill /4}

\[
\begin{aligned}
4x-1&=6\\
4x&=6+1\\
4x&=7\\
x&=\dfrac{7}{4}.
\end{aligned}
\]

L'ensemble des solutions est \(\textcolor{blue}{S=\left\{\dfrac{7}{4}\right\}}\).

\axegraduesolution{1.75}{x=\frac{7}{4}}

\textbf{Exercice 4 -- Intervalles (bonus) \hfill /1}

La borne \(-5\) est exclue (\(<\)) et la borne \(1\) est incluse (\(\leqslant\)), donc l'intervalle recherché est
\[
\textcolor{blue}{\left]-5\ ;\ 1\right]}.
\]

\newpage

% =========================================================
% PAGE 3 : VERSION C
% =========================================================

\interroversion{C}

\textbf{Exercice 1 -- Écriture décimale de fractions \hfill /2}

\begin{enumerate}
\item \(\dfrac{13}{5}=\textcolor{blue}{2{,}6}\)

\item \(-\dfrac{7}{4}=\textcolor{blue}{-1{,}75}\)
\end{enumerate}

\textbf{Exercice 2 -- Solution d'une équation \hfill /4}

Testons \(x=2\) dans \((E_1)\) :
\[
5\times2-3=10-3=7.
\]
On retrouve bien le membre de droite, donc \(\textcolor{blue}{2 \text{ est solution de } (E_1)}\).

Testons \(x=2\) dans \((E_2)\) :
\[
4\times2+1=8+1=9\neq17.
\]
Donc \(\textcolor{blue}{2 \text{ n'est pas solution de } (E_2)}\).

\textbf{Exercice 3 -- Résolution d'équation \hfill /4}

\[
\begin{aligned}
2x+5&=2\\
2x&=2-5\\
2x&=-3\\
x&=-\dfrac{3}{2}.
\end{aligned}
\]

L'ensemble des solutions est \(\textcolor{blue}{S=\left\{-\dfrac{3}{2}\right\}}\).

\axegraduesolution{-1.5}{x=-\frac{3}{2}}

\textbf{Exercice 4 -- Intervalles (bonus) \hfill /1}

La borne \(-1\) est incluse (\(\leqslant\)) et la borne \(3\) est exclue (\(<\)), donc l'intervalle recherché est
\[
\textcolor{blue}{\left[-1\ ;\ 3\right[}.
\]

\newpage

% =========================================================
% PAGE 4 : VERSION D
% =========================================================

\interroversion{D}

\textbf{Exercice 1 -- Écriture décimale de fractions \hfill /2}

\begin{enumerate}
\item \(\dfrac{3}{2}=\textcolor{blue}{1{,}5}\)

\item \(-\dfrac{17}{5}=\textcolor{blue}{-3{,}4}\)
\end{enumerate}

\textbf{Exercice 2 -- Solution d'une équation \hfill /4}

Testons \(x=5\) dans \((E_1)\) :
\[
2\times5-3=10-3=7.
\]
On retrouve bien le membre de droite, donc \(\textcolor{blue}{5 \text{ est solution de } (E_1)}\).

Testons \(x=5\) dans \((E_2)\) :
\[
3\times5+2=15+2=17\neq20.
\]
Donc \(\textcolor{blue}{5 \text{ n'est pas solution de } (E_2)}\).

\textbf{Exercice 3 -- Résolution d'équation \hfill /4}

\[
\begin{aligned}
4x+3&=-2\\
4x&=-2-3\\
4x&=-5\\
x&=-\dfrac{5}{4}.
\end{aligned}
\]

L'ensemble des solutions est \(\textcolor{blue}{S=\left\{-\dfrac{5}{4}\right\}}\).

\axegraduesolution{-1.25}{x=-\frac{5}{4}}

\textbf{Exercice 4 -- Intervalles (bonus) \hfill /1}

La borne \(-4\) est exclue (\(<\)) et la borne \(2\) est incluse (\(\leqslant\)), donc l'intervalle recherché est
\[
\textcolor{blue}{\left]-4\ ;\ 2\right]}.
\]

\end{document}
